\clearpage
\SourceLeaf{189}
\thispagestyle{empty}
\pdfbookmark[0]{Besprechungen fremder Schriften}{other-notices-title}
\begin{center}
\vspace*{1in}{\Large BESPRECHUNGEN\par}
\vspace{.4in}{\large FREMDER\par}
\vspace{.4in}{\Large SCHRIFTEN.\par}
\end{center}

\clearpage
\SourceLeaf{191}
\renewcommand{\BookRunningTitle}{Besprechungen fremder Schriften}
\pdfbookmark[1]{Recherches sur l'irréductibilité. Anzeige 1809, 11. März}{notice-1809-03-11}
\begin{center}
\rule{.86\linewidth}{.3pt}\par
Göttingische gelehrte Anzeigen.\quad 1809, 11. März.\par
\rule{.86\linewidth}{.3pt}
\end{center}
\begin{center}
\emph{Recherches sur l'irréductibilité Arithmétique et Géométrique des nombres et de leurs puissances.}\quad 1808.\quad (Ohne Druckort.\quad 25 S. in großem Quartformat.)
\end{center}

Eine Schrift mit dem Ziel, irrationale Wurzelgrößen in der Form rationaler Größen darzustellen. Wir müssen uns damit begnügen, die Freunde der Mathematik auf dieses kleine Werk aufmerksam gemacht zu haben. Denn der Umfang dieser Blätter erlaubt uns nicht, hier ausführlicher den eigentümlichen Blickwinkel des Verfassers und seine von der üblichen ganz abweichende Behandlung der Wurzelgrößen darzustellen und zu prüfen.

\Needspace{8\baselineskip}
\pdfbookmark[1]{Chernac, Cribrum Arithmeticum. Anzeige 1812, 23. März}{notice-1812-03-23}
\begin{center}
\rule{.86\linewidth}{.3pt}\par
Göttingische Gelehrte Anzeigen.\quad 1812, 23. März.\par
\rule{.86\linewidth}{.3pt}
\end{center}
\emph{Cribrum Arithmeticum, sive tabula continens numeros primos a compositis segregatos, occurrentes in serie numerorum ab unitate progredientium usque ad decies centena milia et ultra haec ad viginti millia} $(1020000)$. \emph{Numeris compositis, per} $2$, $3$, $5$ \emph{non dividuis, adscripti sunt divisores simplices, non minimi tantum, sed omnino}
\SourceLeaf{192}
\emph{omnes. Confecit} \textsc{Ladislaus Chernac}, \emph{Pannonius, A. L. M. Philos. et Medic. Doctor, in almo lyceo Daventriensi philosophiae professor. Daventriae} 1811. (Auf Kosten des Verfassers, gedruckt bei J. H. Lange.\quad XXII und 1022 S. in großem Quartformat.)

Der vollständige Titel dieses wichtigen und sehr verdienstvollen Werks bezeichnet den Inhalt bereits hinreichend: Es ist eine durch ebenso sorgfältige wie mühsame mehrjährige Arbeit berechnete Tafel aller Primfaktoren sämtlicher durch $2$, $3$ und $5$ nicht teilbaren Zahlen von $1$ bis $1020000$, sauber und, soweit wir bei gelegentlichen Prüfungen feststellen konnten, sehr korrekt gedruckt. Wie wertvoll ein solches Geschenk an die Arithmetik ist, kann jeder leicht beurteilen, der häufig mit größeren Zahlen rechnet. Der Verfasser verdient doppelten Dank: sowohl für die höchst mühsame Arbeit selbst, durch die er seinen Namen den unvergesslichen Namen von \textsc{Rhaeticus}, \textsc{Pitiscus}, \textsc{Brigg}, \textsc{Vlacq}, \textsc{Wolfram}, \textsc{Taylor} und anderen zugesellt hat, als auch für den gewiss sehr erheblichen Druckaufwand, für den sich sonst wohl kaum ein Verleger gefunden hätte. Schon öfter wurden solche Tafeln berechnet, wenn auch meist in geringerem Umfang, blieben aber entweder ganz im Manuskript oder wurden nicht vollständig gedruckt. Bekanntlich ermunterte \textsc{Lambert} früher nach besten Kräften zur Fortsetzung der oft gedruckten \textsc{Pell}schen Tafel, die bis $100000$ reicht. Nach einer Mitteilung von \textsc{Bernoulli} in \textsc{Lambert}s Briefwechsel hatte \textsc{Oberreit} sie bis $500000$ fortgeführt; eine Abschrift davon war in \textsc{Schulze}s Hände gelangt. \textsc{Anton Felkel} hatte sie, wie in der Monatl. Correspondenz, Band 2, S. 223 berichtet wird, im Manuskript bis zu zwei Millionen vollendet und wollte sie später bis $2460000$ liefern. Was davon in Wien auf öffentliche Kosten bereits gedruckt worden war, wurde jedoch im Türkenkrieg zu Patronen verarbeitet, weil sich keine Käufer fanden! So ging eine verdienstvolle Arbeit vieler Jahre für die Öffentlichkeit verloren. Umso mehr hielten wir es für unsere Pflicht, hier das Erscheinen des gegenwärtigen Werks anzuzeigen. Die erste Million steht nun jedem zum Gebrauch zur Verfügung; wer Gelegenheit und Eifer für diesen Gegenstand hat, möge daher seine Mühe auf das Weitere richten.

\SourceLeaf{193}
\Needspace{8\baselineskip}
\pdfbookmark[1]{Burckhardt, deuxième million. Anzeige 1814, 3. November}{notice-1814-11-03}
\begin{center}
\rule{.86\linewidth}{.3pt}\par
Göttingische gelehrte Anzeigen.\quad 1814, 3. November.\par
\rule{.86\linewidth}{.3pt}
\end{center}
\emph{Tables des diviseurs pour tous les nombres du deuxième million, ou plus exactement depuis} $1020000$ \emph{à} $2028000$, \emph{avec les nombres premiers qui s'y trouvent. Par} \textsc{J. Ch. Burckhardt}, \emph{membre de l'institut impérial, du bureau des longitudes de France, et de plusieurs autres sociétés savantes. Paris}, 1814. \emph{M\textsuperscript{me} V\textsuperscript{e} Courcier}. (VIII. und 112 S. in Folio.)

Früher, als wir bei der Anzeige der \textsc{Chernac}schen Faktorentafel für die erste Million zu hoffen gewagt hätten, können wir bereits die Vollendung und das Erscheinen einer ähnlichen Tafel für die zweite Million melden. Der verdienstvolle Verfasser, dessen Name allein schon größte Sorgfalt und Genauigkeit verbürgt, hat durch diese mühsame Arbeit alle Freunde der Arithmetik sehr verpflichtet. \textsc{Chernac}s Tafel für die erste Million enthält alle Primfaktoren, \textsc{Burckhardt}s Tafel für die zweite dagegen jeweils nur den kleinsten Teiler. Die vollständige Zerlegung einer Zahl der zweiten Million erfordert daher die Division durch den kleinsten Teiler und das Nachschlagen des Quotienten in der \textsc{Chernac}schen Tafel. Diese kleine Mühe fällt jedoch gegenüber dem großen Vorteil, die Tafel auf so viel kleinerem Raum zu besitzen, gar nicht ins Gewicht. Dabei bleibt die Aussicht, die Tafel mit der Zeit noch bis zu zehn Millionen erweitert zu sehen. Die Unterbringung in dem kleinen Band ermöglichte der Verfasser teils durch die Beschränkung auf den kleinsten Teiler, teils durch einen möglichst sparsamen Druck. Bezeichnet $a$ unbestimmt eine der achtzig Zahlen unter $300$, die durch $2$, $3$ und $5$ nicht teilbar sind, so ist jede durch $2$, $3$ und $5$ nicht teilbare Zahl in der Form $300n+a$ enthalten. Alle achtzig Zahlen, für die $n$ denselben Wert hat, stehen in einer einzigen senkrechten Spalte; jede Seite enthält dreißig solcher Spalten. Jede Seite umfasst also von neuntausend in natürlicher Reihenfolge aufeinanderfolgenden Zahlen alle, die durch $2$, $3$ oder $5$ nicht teilbar sind.

\Needspace{6\baselineskip}
Die Methode, mit der Herr \textsc{Burckhardt} seine Tafel erstellt hat, verdient hier besondere Erwähnung. Er ließ ein Netz in Kupfer stechen, in dem $81$ waagerechte und $78$ senkrechte Linien ein in $80\times77$, also $6160$ kleine Quadrate geteiltes Rechteck bildeten, und davon die nötige Anzahl Abzüge herstellen. An der Seite konnten die achtzig Werte
\SourceLeaf{194}
von $a$ gleich mit eingestochen werden; die Werte von $300n$ schrieb man in fortlaufender Reihenfolge mit der Feder über die $77$ senkrechten Spalten. Jedes Blatt stellt so alle durch $2$, $3$ und $5$ nicht teilbaren Zahlen dar, die sich unter jeweils $23100$ in natürlicher Reihenfolge aufeinanderfolgenden Zahlen befinden; $44$ Blätter genügen, um eine ganze Million zu umfassen. Man sieht leicht, dass die Zahlen, deren kleinster Teiler $7$ oder $11$ ist, auf jedem folgenden Blatt in derselben Reihenfolge wiederkehren. Diese Teiler konnten daher gleich in die Kupferplatte gestochen werden und erschienen so auf jedem Blatt von selbst an den richtigen Stellen. Um nun die folgenden Teiler, beispielsweise $13$, einzutragen, nahm Herr B. von einem überzähligen Blatt der Breite nach nur $13$ Spalten. Er betrachtete dieses als den Anfang seiner Tafel und schnitt alle Quadrate aus, die den Teiler $13$ enthalten mussten. Dieses Gitter brauchte er dann nur auf die ersten dreizehn Spalten des ersten Blattes, danach auf die nächsten dreizehn usw. zu legen, um sofort alle Stellen zu sehen, die, soweit sie nicht schon $7$ oder $11$ enthielten, mit $13$ auszufüllen waren. Ebenso verfuhr er danach mit dem Teiler $17$ usw. Bis zum Teiler $73$ reichten die überzähligen Blätter auf diese Weise aus. Für die größeren Teiler $79$, $83$ usw. scheint Herr B. den Rahmen aus zwei oder mehreren Teilen zusammengesetzt zu haben. Bei den Teilern über $500$ zog Herr B. dagegen vor, die Vielfachen durch Addition zu ermitteln, wobei er für den anderen Faktor nur Primzahlen zu nehmen brauchte. Wir halten das ganze Verfahren für höchst zweckmäßig und würden es allen zur Nachahmung empfehlen, die die Tafel vielleicht noch weiter fortsetzen möchten. Für die dritte und vierte Million hat der Verfasser inzwischen selbst schon einen großen Teil der Rechnungen ausgeführt. Wir haben daher begründete Hoffnung, auch diese bald im Druck zu sehen.

\begingroup\setlength{\topsep}{0pt}\setlength{\partopsep}{0pt}
\Needspace{22\baselineskip}
\pdfbookmark[1]{Burckhardt, troisième million. Anzeige 1816, 7. November}{notice-1816-11-07}
\begin{center}
\rule{.86\linewidth}{.3pt}\par
Göttingische gelehrte Anzeigen.\quad 1816, 7. November.\par
\rule{.86\linewidth}{.3pt}
\end{center}
\emph{Tables des diviseurs pour tous les nombres du troisième million, ou plus exactement, depuis} $2028000$ \emph{à} $3036000$, \emph{avec les nombres premiers qui s'y trouvent; par} \textsc{J. Chr. Burckhardt}, \emph{membre de l'académie royale des sciences, du bureau des }
\SourceLeaf{195}
\emph{longitudes de France et de plusieurs autres sociétés savantes. Paris} 1816. \emph{M\textsuperscript{me} V\textsuperscript{e} Courcier}. (112 S. in Folio.)

Da wir bei der Anzeige der Faktorentafel für die zweite Million bereits die Berechnungsmethode des verdienstvollen Verfassers und die Einrichtung der Tafel selbst ausführlich beschrieben haben, können wir uns hier mit der bloßen Anzeige des Erscheinens der Tafel für die dritte Million begnügen. Bald dürfen wir vom Verfasser auch noch die auf dieselbe Art dargestellte Tafel für die \emph{erste} Million erwarten, sodass die ganze Tafel bis über drei Millionen nur einen mäßig großen Band ausmachen wird. Dafür gebührt dem Verfasser der Dank aller Freunde der Arithmetik. Sie sehen durch diese mühsame Arbeit ein Bedürfnis in einem Umfang erfüllt, der alles weit übersteigt, was man noch vor wenigen Jahren zu hoffen wagen konnte.

\Needspace{8\baselineskip}
\pdfbookmark[1]{Burckhardt, premier million. Anzeige 1817, 9. August}{notice-1817-08-09}
\begin{center}
\rule{.86\linewidth}{.3pt}\par
Göttingische gelehrte Anzeigen.\quad 1817, 9. August.\par
\rule{.86\linewidth}{.3pt}
\end{center}
\emph{Tables des diviseurs, pour tous les nombres du premier million, ou plus exactement depuis} $1$ \emph{à} $1020000$, \emph{avec les nombres premiers qui s'y trouvent; par} \textsc{J. Chr. Burckhardt}, \emph{membre de l'académie des sciences dans l'institut royal, du bureau des longitudes de France, et de plusieurs autres sociétés savantes. Paris} 1817. \emph{M\textsuperscript{me} V\textsuperscript{e} Courcier}. (114 S. in Folio.)

Mit Verweis auf unsere Anzeigen der Tafeln für die zweite und dritte Million kündigen wir jetzt nur das tatsächliche\textperiodcentered\ Erscheinen dieser Faktorentafeln für die erste Million an. Wir besitzen damit nun ein zusammenhängendes Ganzes für die ersten drei Millionen. Für die gegenwärtige erste Million benutzte der Verfasser teils das \emph{Cribrum Arithmeticum} von \textsc{Chernac}, teils eine handschriftliche Tafel von \textsc{Schenmark}, die die Bibliothek des Königlichen Instituts besitzt. Letztere war allerdings nicht ganz mit aller wünschenswerten Sorgfalt erstellt worden. Wenn beide voneinander abwichen, war die Entscheidung, welche von beiden recht hatte, oft ziemlich mühsam. In der \textsc{Chernac}schen Tafel fand sich nur eine sehr geringe Anzahl von Fehlern, die Herr \textsc{Burckhardt} hier mitgeteilt hat.

\SourceLeaf{196}
Auch für die vierte, fünfte und sechste Million hat der Verfasser die Materialien bereits größtenteils bereit. Er bietet an, diese Fortsetzung zu liefern, wenn der Verleger durch einen hinreichenden Absatz der ersten drei Millionen ermutigt wird. Es wäre in der Tat sehr zu bedauern, wenn die Früchte einer so mühsamen und nützlichen Arbeit der Welt vorenthalten würden.
\endgroup
